\section{Results}
\label{sec:results}

\subsection{Headline comparison}

Table~\ref{tab:main} reports the matched-step comparison at step $11{,}500$,
where all fifteen runs are present.

\begin{table}[t]
\centering
\begin{tabular}{lrrrc}
\toprule
Algorithm & Reward & Content & Sentiment & Seeds \\
\midrule
R-REBEL ($\ell_1$-std) & 39.34 & 56.25 & 53.70 & 3/3 \\
R-REBEL (Huber-std) & 39.27 & 58.62 & 49.55 & 3/3 \\
R-REBEL ($\ell_1$-ent) & 37.49 & 58.17 & 48.15 & 3/3 \\
GRPO + entropy & 33.36 & 53.57 & 50.57 & 3/3 \\
GRPO & 30.19 & 55.92 & 43.03 & 3/3 \\
\bottomrule
\end{tabular}

\caption{Mean reward, content, and sentiment at matched step $11{,}500$ ---
the largest step at which every arm has all three seeds. Reward is averaged
over the evaluation set and then over the $\lmb$ grid $\{0,\dots,0.9\}$. All
quantities are on a $0$--$100$ scale. Note the evaluation set here is the
$10$-sentence development split, not the held-out test split
(Section~\ref{sec:setup-exp}).}
\label{tab:main}
\end{table}

Three observations, in decreasing order of robustness.

\paragraph{1. The family gap is large and unambiguous.}
The best \rrebel{} variant scores $39.34$ against $33.36$ for the best GRPO arm:
a gap of $\mathbf{5.98}$ points, or six times the $1.0$-point noise floor. The
separation survives at the level of individual runs, which is the strongest form
the claim can take: the \emph{worst} \rrebel{} seed ($36.82$) beats the
\emph{best} GRPO seed ($34.65$) by $2.16$ points. There is no overlap between
the two families' seed distributions, so the ordering does not depend on seed
selection or on averaging. Averaged over the second half of training
(steps $\ge 6000$), the same picture holds: $39.37$ and $39.44$ for
$\ell_1$-std and Huber-std versus $31.11$ and $33.17$ for the two GRPO arms.

\paragraph{2. The top two \rrebel{} variants are tied.}
$\ell_1$-std ($39.34$) and Huber-std ($39.27$) differ by $0.07$ points, far
inside the noise floor, and their per-seed ranges overlap almost completely
($39.0$--$39.9$ versus $38.7$--$40.0$). We therefore report them as tied and
draw no conclusion about which robust discrepancy is preferable.
$\ell_1$-ent, at $37.49$, is a genuine $1.85$ points behind --- adding an
entropy bonus to $\ell_1$-std \emph{hurt}, which is consistent with the claim in
Section~\ref{sec:rrebel} that \rrebel{} does not need an explicit exploration
term because its fixed point is already a distribution.

\paragraph{3. GRPO's deficit is in sentiment, not content.}
The component columns of Table~\ref{tab:main} locate the failure.
Plain GRPO's content score ($55.92$) is competitive with $\ell_1$-std's
($56.25$), but its sentiment score is $43.03$ against $53.70$ --- a $10.7$-point
shortfall that accounts for essentially the entire reward gap. Adding an entropy
bonus moves GRPO's sentiment up to $50.57$ but at the cost of content, which
falls to $53.57$, the lowest of any arm. In other words the GRPO arms can
purchase one objective or the other but cannot hold both at once, which is
precisely the pathology a $\lmb$-conditioned policy must avoid.
Figure~\ref{fig:components} shows this as a training-time trajectory: the
content panels are tightly bunched across all five arms, while the sentiment
panel separates the families cleanly.

\begin{figure}[t]
\centering
\includegraphics[width=0.82\textwidth]{learning_curves.pdf}
\caption{Mean reward against training step, averaged over three seeds, with the
band spanning the seed minimum to maximum. The two families separate by step
$\sim\!2000$ and never re-cross. Blue denotes \rrebel{} variants, orange the
GRPO arms.}
\label{fig:learning}
\end{figure}

\begin{figure}[t]
\centering
\includegraphics[width=\textwidth]{components.pdf}
\caption{The reward gap decomposed. Content (left) is comparable across all
arms; sentiment (right) is where plain GRPO falls away. Both panels share a
$y$-scale, so the difference in spread between them is directly comparable.}
\label{fig:components}
\end{figure}

\subsection{Learning dynamics}

Figure~\ref{fig:learning} shows that the outcome is not a late-training
artifact. All arms improve rapidly over the first $\sim\!1500$ steps; the two
families have separated by step $2000$ and never cross again. The \rrebel{}
variants climb to roughly $39$--$40$ and then plateau, whereas both GRPO arms
flatten early --- plain GRPO at $\sim\!31$ and GRPO-plus-entropy at $\sim\!33$
--- and stay flat for the remaining $10{,}000$ steps. The seed bands are also
informative: GRPO's band is markedly wider (its seeds at step $11{,}500$ span
$26.33$ to $34.65$, an $8.3$-point spread, versus $36.82$--$40.00$, a
$3.2$-point spread, across all nine \rrebel{} seeds), so the baseline is both
worse on average and less reliable.

\paragraph{Plain GRPO does not plateau --- it never learns at all.}
This distinction matters, and reading Figure~\ref{fig:learning} as ``GRPO
converges lower'' would be too generous to our own result. Measured from its
\emph{first} evaluation at step $500$ to step $11{,}500$, the plain GRPO arm
moves from $30.38$ to $30.19$: a change of $\mathbf{-0.18}$ points over
$11{,}000$ training steps, inside the noise floor and slightly negative. The
per-seed deltas are $+0.11$, $+1.09$ and $-1.75$. For comparison, over the same
interval \rrebel{} ($\ell_1$-std) gains $+6.74$, $+6.23$ and $+7.77$, and
GRPO-plus-entropy gains $+17.74$, $+3.23$ and $+16.77$. Two corroborating
diagnostics from the training logs: the logged KL sits at $9.824$ against the
$k_3$ estimator's analytic ceiling of $9.8249$ for a $50{,}257$-token
vocabulary --- attained only when the sampled token has probability $\approx 1$
--- and the count of distinct prompt tokens explored per training cycle has a
median of $5$, which for a $5$-token prompt means a single frozen,
$\lmb$-independent prompt.

So the plain GRPO number in Table~\ref{tab:main} is the score of a policy that
collapsed by step $500$ and never moved again. It is a \emph{diagnostic}, not a
converged baseline, and we do not present it as one: \textbf{the defensible GRPO
baseline in this paper is the GRPO-plus-entropy arm} at $33.36$, which does
learn. The headline comparison should be read against that arm --- a $5.98$-point
gap --- and plain GRPO read as evidence for the collapse mechanism of
Section~\ref{sec:grpo} rather than as a tuned competitor.
\todo{re-run plain GRPO with a competent (non-uniform) KL reference, per the
TODO in Section~\ref{sec:algos-caveat}, before claiming anything about GRPO's
ceiling as an algorithm.}

\subsection{The baseline was debugged and tuned, and still trails}
\label{sec:baseline-defense}

The most likely objection to a result of this shape is that the baseline was
misconfigured. It was, initially --- and fixing it is part of the result. We
report the full sequence in Figure~\ref{fig:baseline}, all at matched step
$4500$:

\begin{center}
\begin{tabular}{lrl}
\toprule
GRPO configuration & Mean reward & Change \\
\midrule
Rolling reference anchor (defective) & $27.16$ & --- \\
Pinned reference anchor (Section~\ref{sec:grpo}) & $31.30$ & $+4.14$ \\
\quad + entropy bonus $\alpha_H = 0.01$ & $33.10$ & $+1.80$ \\
\midrule
\rrebel{} (all variants, same step) & $39.26$ & $+6.16$ remaining \\
\bottomrule
\end{tabular}
\end{center}

The first row is the defect described in Section~\ref{sec:grpo}: the KL
reference was being re-copied from the live policy every $150$ steps, roughly
eighty times over a run, which zeroed the penalty and its gradient at every
refresh. Under that configuration GRPO's reward curve was flat from step $500$
onward, consistent with the entropy collapse that Section~\ref{sec:grpo}
predicts when the only regularizer is inert. Pinning the anchor is worth
$+4.14$ points. Giving GRPO the same entropy-bonus lever that $\ell_1$-ent has
is worth a further $+1.80$.

Both interventions behave as designed, which is evidence that they were
implemented correctly rather than merely tried. The entropy bonus in particular
starts \emph{worse} than plain GRPO and only overtakes it from step
$\sim\!1500$ onward --- the signature of a term that buys exploration at the
cost of early exploitation.

We must, however, retract one part of the defense. It would be convenient to say
that the two GRPO arms are two \emph{independent} chances for the baseline, so
that their agreeing on a $\sim\!31$--$33$ ceiling is strong evidence about GRPO
as an algorithm. That claim does not survive scrutiny. As
Section~\ref{sec:algos-caveat} shows, the pinned reference is the uniform
distribution, so plain GRPO already carries an implicit entropy coefficient of
$0.04$; the ``GRPO + entropy'' arm merely raises the total to $\approx\!0.05$.
The two arms are one configuration at two nearby settings of the same
anti-collapse lever, which is a natural explanation for their landing in the
same place, and it is \emph{not} independent corroboration. What the table above
does support is narrower but still real: a diagnosed defect was fixed and was
worth $+4.14$, an additional exploration lever was worth $+1.80$, and a
$6$-point gap to \rrebel{} remained after both.

\todo{the direct collapse evidence --- per-checkpoint counts of distinct emitted
prompts, which during the campaign read $2$--$19$ for defective GRPO against
$\sim\!97$ for \rrebel{} --- came from the per-$\lmb$ evaluation files and did
not survive the purge (Section~\ref{sec:provenance}). It is deliberately not
quoted as a measured result here. The surviving tracker metric
(\texttt{num\_tokens\_explored}) cannot substitute: it is a per-process
cumulative set that restarts on every $4$-hour resubmit, so its peak is
dominated by early-training exploration and does not measure steady-state
collapse. The rerun regenerates the original statistic; restore this paragraph's
numbers from it.}

\begin{figure}[t]
\centering
\includegraphics[width=0.82\textwidth]{baseline_debug.pdf}
\caption{The baseline-defense figure. The defective rolling-anchor GRPO (grey)
is flat at $\sim\!27$; pinning the anchor lifts it to $\sim\!31$ (solid orange);
adding an entropy bonus reaches $\sim\!33$ (dashed orange). The blue band is the
envelope over all three \rrebel{} variants and all their seeds. The rolling-anchor
runs were terminated once the defect was identified, hence their shorter extent.}
\label{fig:baseline}
\end{figure}

\subsection{The content--sentiment tradeoff curve}
\label{sec:tradeoff}

\begin{figure}[t]
\centering
\includegraphics[width=0.6\textwidth]{tradeoff.pdf}
\caption{Content against sentiment at step $11{,}500$. Hollow markers are
individual seeds, filled markers each arm's mean. \textbf{These are
$\lmb$-averaged operating points, not a frontier}: each marker collapses the
whole $\lmb$ grid to one point. The per-$\lmb$ curve is pending
(Section~\ref{sec:provenance}). Even so, the \rrebel{} points sit up and to the
right of the GRPO points, i.e.\ better on both objectives simultaneously.}
\label{fig:tradeoff}
\end{figure}

The intended headline figure of this paper is the \emph{per-$\lmb$} curve: for
each arm, the locus of (content, sentiment) pairs obtained as $\lmb$ sweeps
$0 \to 0.9$, which shows not just that \rrebel{} scores higher but that it
commands a strictly better frontier. That figure cannot yet be drawn --- the
per-$\lmb$ evaluation arrays for the completed campaign were lost, as
Section~\ref{sec:provenance} details, and regenerating them requires
checkpoints, which are being retrained.

Figure~\ref{fig:tradeoff} therefore shows what the surviving data does support:
each run's single $\lmb$-averaged operating point. This is deliberately
\emph{not} presented as a frontier. It still carries a real claim --- the
\rrebel{} points lie up and to the right of the GRPO points, so \rrebel{}
dominates on both objectives at once rather than trading one for the other ---
but the shape of each arm's curve, and in particular whether GRPO's curve is
uniformly dominated or only dominated in the interior of the $\lmb$ range,
remains open until the rerun completes.
\todo{replace Figure~\ref{fig:tradeoff} with the true per-$\lmb$ frontier once
\texttt{slurm/eval\_v4.slurm} has produced
\texttt{results/v4/*/test/output*.json}; \texttt{analysis/paper\_results.py}
switches the figure over automatically and reports which version it drew.}


\subsection{Results at step 12{,}000}

\begin{table}[t]
\centering
\begin{tabular}{lrrrc}
\toprule
Algorithm & Reward & Content & Sentiment & Seeds \\
\midrule
R-REBEL ($\ell_1$-std) & 39.47 & 55.87 & 53.63 & 3/3 \\
R-REBEL (Huber-std) & 39.30 & 60.34 & 47.78 & 2/3 \\
R-REBEL ($\ell_1$-ent) & 37.80 & 58.88 & 48.72 & 2/3 \\
GRPO + entropy & 32.68 & 54.12 & 46.93 & 2/3 \\
GRPO & 30.91 & 55.24 & 45.09 & 3/3 \\
\bottomrule
\end{tabular}

\caption{The same quantities at step $12{,}000$, the nominal end of training.
Three runs have no evaluation at this step (Section~\ref{sec:provenance}), so
the seed counts differ by arm; the ordering and the size of the family gap are
unchanged from Table~\ref{tab:main}.}
\label{tab:12000}
\end{table}

Table~\ref{tab:12000} reports the nominal end of training. The conclusions do
not change: $\ell_1$-std ($39.47$) and Huber-std ($39.30$) remain tied within
noise, $\ell_1$-ent trails at $37.80$, and the GRPO arms sit at $32.68$ and
$30.91$. Because three runs are missing their $12{,}000$ evaluation, this table
is seed-unbalanced across arms and Table~\ref{tab:main} is the comparison we
rely on.

\subsection{Data provenance, and what is missing}
\label{sec:provenance}

Reporting this honestly requires a note on provenance, because it also bounds
what can currently be reproduced.

The campaign ran to completion, but its artifacts were subsequently destroyed:
the scratch filesystem holding all checkpoints and all per-$\lmb$ evaluation
files was purged, leaving zero checkpoints and zero evaluation files. Every
number in this paper was therefore recovered from the experiment tracker's
retained scalar history, by a reconstruction that is validated three ways per
run and refuses to emit data if any check fails: surviving true-step markers
must agree with the reconstructed step, the distinct-step count must be
consistent with the final step, and the reconstructed step grid must be
perfectly regular. All fifteen runs pass. Two consequences:

\begin{enumerate}\itemsep2pt
  \item \textbf{The per-$\lmb$ arrays are gone and are not recoverable from the
        tracker}, because the evaluation routine averages over the $\lmb$ grid
        \emph{before} logging. This is why Section~\ref{sec:tradeoff} shows
        operating points rather than a frontier. One seed per arm is being
        retrained specifically to regenerate them.
  \item \textbf{Three runs lack a step-$12{,}000$ evaluation.} Their final
        training cycle ran while the tracker was unreachable; the trainer's
        fallback keeps training without logging, so that cycle's evaluations
        were written only to the since-purged disk copies. This is why step
        $11{,}500$, not $12{,}000$, is the balanced comparison point.
\end{enumerate}

Two process changes follow, and are already in place: evaluation artifacts are
now copied off the purgeable filesystem after every training cycle, and the
recovery-and-validation script is committed alongside the exported data so that
every table in this paper can be regenerated from a single command.
